\documentclass[11pt,a4paper]{article}
\usepackage[T1]{fontenc}
\usepackage[utf8]{inputenc}
\usepackage{lmodern}
\usepackage[margin=2cm]{geometry}
\usepackage{enumitem}
\usepackage{xcolor}
\usepackage[hidelinks]{hyperref}
\usepackage{titlesec}
\usepackage{needspace}

\definecolor{accent}{RGB}{30,60,110}
\pagestyle{empty}
\setlength{\parindent}{0pt}
\titleformat{\section}{\Needspace{6\baselineskip}\large\bfseries\color{accent}}{}{0pt}{}[\vspace{-0.6em}\color{accent}\rule{\textwidth}{0.4pt}]
\titlespacing*{\section}{0pt}{1.0em}{0.6em}
\setlist[itemize]{leftmargin=1.2em,itemsep=0.15em,topsep=0.2em}

% \entry{dates}{title}{place}{details}
\newcommand{\entry}[4]{%
  \noindent\makebox[2.6cm][l]{#1}%
  \parbox[t]{\dimexpr\textwidth-2.6cm\relax}{\textbf{#2}%
    \if\relax\detokenize{#3}\relax\else, #3\fi
    \if\relax\detokenize{#4}\relax\else\\#4\fi}\par\vspace{0.45em}}
% Visible marker for details still to confirm. Remove before sending.
\newcommand{\todo}[1]{\textcolor{red}{[#1]}}

\begin{document}

{\LARGE\bfseries Aleksandras Melnik}\\[0.3em]
Logic \textperiodcentered{} Philosophy of Mathematics and Physics \textperiodcentered{} Category Theory\\[0.3em]
\href{mailto:aaamelnik@gmail.com}{aaamelnik@gmail.com} \quad\textperiodcentered\quad
\href{https://sketch-of-an-elephant.github.io/aleksandras-melnik}{sketch-of-an-elephant.github.io/aleksandras-melnik}\\
\href{https://github.com/sketch-of-an-elephant}{github.com/sketch-of-an-elephant}

\section{Research interests}
Equivalence of scientific theories (definitional, Morita and categorical equivalence);
model theory and category theory applied to physical theories;
symmetry and structure in general relativity and quantum theory;
categorical logic and the foundations of mathematics.

\section{Education}
\entry{2024 -- 2027}{MA in Logic and Philosophy of Science}{Munich Center for Mathematical Philosophy (MCMP), LMU Munich}
  {Expected April 2027. DAAD scholarship for master's studies.\\
   Thesis supervisor: John Dougherty.\\
   Uses logic and category theory to ask when two formulations of a physical theory are equivalent.}
\entry{2020 -- 2024}{BSc in Mathematics and Applications of Mathematics}{Vilnius University}{}
\entry{2016 -- 2020}{Vilnius Lyceum}{Lithuania}{100\% in the national mathematics exam.}

\section{Honours and awards}
\entry{2024}{DAAD Scholarship}{German Academic Exchange Service}{Full scholarship for master's studies in Germany.}
\entry{2020}{Bronze medal}{International Mathematical Olympiad (IMO)}{}

\section{Talks}
\entry{Jan 2026}{Grounding and Explanation Converging in Mathematics}{Nineteenth Annual Cambridge Graduate Conference in Philosophy of Mathematics and Logic, University of Cambridge}{}
\entry{Nov 2025}{Justifying the Laws of Logic}{Vilnius Philosophy Graduate Conference, Vilnius University}{}
\entry{}{Category-theoretic models for mereology}{Vilnius University Philosophy Student Conference}{}

\section{Research experience}
\entry{2022}{Visiting Researcher}{Institute of Computer Science, Czech Academy of Sciences}{}
\entry{2022}{Research Assistant}{Institute of Mathematics, Vilnius University}{}

\section{Research seminars (Vilnius University)}
Participant in weekly research seminars, and curator (organiser) of three:
\begin{itemize}
  \item \textbf{Curator:} Algebraic Topology (2022--23), Intuitionistic Logic (2023), Mathematical Quantum Mechanics (2024)
  \item Category Theory (2022--24); Mathematical Logic (2021--22); Axiomatic Theories of Truth (2022)
  \item Extremal Combinatorics (2019--24); Algebraic Geometry (2023--24); Algebraic Number Theory (2022);
        Dynamical Systems (2022--24); Network Theory (2023--24)
\end{itemize}

\section{Teaching}
\entry{2021 -- 2024}{Lecturer, Lithuanian national olympiad training camps}{}
  {Training the national teams for the IMO (2021--24), the Middle European Mathematical Olympiad (2021--23) and the Baltic Way competition (2021--23).}

\section{Competitions}
\begin{itemize}
  \item International Mathematical Olympiad: bronze medal, 2020
  \item Vojtěch Jarník International Mathematical Competition: Lithuanian team, 2024
  \item Baltic Way: Lithuanian team, 2018 and 2019
  \item Middle European Mathematical Olympiad: Lithuanian team, 2016 and 2017
  \item International Philosophy Olympiad: Lithuanian team, 2019
\end{itemize}

\section{Other experience}
\entry{2023 -- 2024}{Data Engineer and Analyst}{Adtractive}{}

\section{Languages}
Lithuanian (native), Ukrainian (native), English (C2, IELTS 8.5), French (fluent), Russian (fluent), German (A1).

\section{Technical skills}
\textbf{Programming:} Python, C++, R, SQL, LookML, \LaTeX.\quad
\textbf{Software:} Mathematica, MATLAB, RStudio, Qiskit, Maxima.

\end{document}
